\section*{Chapter \ref{ch:pl:the-platform-map} --- The Platform Map}

\begin{solution}{exo:pl:the-platform-map:1}
It starts when its last input is ready: the positions, at 20:15. It ends at 20:40. The marks (17:05) and the security master
(19:55) were ready earlier; the positions decide.
\end{solution}

\begin{solution}{exo:pl:the-platform-map:2}
The file's latest safe arrival is 22:35; at 22:50 it is 15 minutes late. The golden copy runs from 22:50 to 23:15, the features
until 01:15, and the backtests end at 07:15: the research report is 15 minutes late. Every other report still has slack (the risk
batch, for example, runs from 23:35 to 05:35).
\end{solution}

\begin{solution}{exo:pl:the-platform-map:3}
Only two datasets derive from it: the reconciliation breaks and the signed P\&L, because it is read only by the reconciliation.
A small \omterm{def:pl:the-platform-map:blast}{blast radius} says how far an error travels through the batch, not how much it matters: a wrong clearing statement can
hide a real break in the firm's positions, which is exactly what the reconciliation exists to find. It needs checking at the
point of use (chapter~22), even though little is built on it.
\end{solution}

\begin{solution}{exo:pl:the-platform-map:4}
The positions are ready 45 minutes after the file (25 minutes of golden copy, 20 of positions), so the batch reads complete
inputs only if the file arrives by $270 - 45 = 225$ minutes after the close, 19:45. The nominal file comes at 19:30: 15 minutes
of margin instead of 45. The earlier start makes the stale read three times more likely to happen on a slightly late night.
\end{solution}

\begin{solution}{exo:pl:the-platform-map:5}
The verified marks are computed from other inputs (broker quotes, consensus data) by another system and carry adjustments; they
are not derived from the library's marks by the map, so they are a different dataset with its own \omterm{def:pl:the-platform-map:sor}{system of record}, product
control. The sign-off reads the verified marks, and reports the difference between the two as a number, because the purpose of
independent verification is that the desk's own marks do not close the books.
\end{solution}

\begin{solution}{exo:pl:the-platform-map:6}
(a) Late information: busts, corrections and late fills arrive after the stream has applied the originals, and a drop copy may
show fills the trading session missed (chapter~17). (b) Reference data: the batch applies the evening's corporate actions from the
security master (a split changes quantities), the stream does not. The books are closed on the batch, reconciled with the
clearing statement; the difference with the stream is investigated as a break.
\end{solution}

\begin{solution}{exo:pl:the-platform-map:7}
Its latest start is $840 - 60 = 780$ minutes after the close, 05:00. No other job's latest start moves: the positions already had
to be ready by 01:00 for the risk batch, well before 05:00. The clearing statement's latest safe arrival becomes 05:00 instead of
05:50, because the forecast now reads it with an earlier deadline than the reconciliation's.
\end{solution}

\begin{solution}{exo:pl:the-platform-map:8}
The risk batch is not on the longest chain. Halved to three hours it ends at 23:15 instead of 02:15, but the batch still ends at
03:55, with the backtests: the longest chain is the corporate-action file, the golden copy, the features and the backtests. And
doubling machines halves a job only if it parallelises perfectly and nothing else limits it (chapters~13 and~20).
\end{solution}

\begin{solution}{pb:pl:the-platform-map:1}
\begin{enumerate}
\item Golden copy 19:30--19:55; positions 19:55--20:15; risk batch 20:15--02:15.
\item The backtests, at 03:55.
\item The corporate-action file, the golden copy, the features and the backtests: $25 + 120 + 360$ minutes from 19:30. The risk
batch's own chain (file, golden copy, positions, risk batch) ends at 02:15, 100 minutes earlier, because the research branch
carries two long jobs.
\item $900 - 715 = 185$ minutes.
\item 285 minutes: its latest start is 01:00 and it starts at 20:15.
\item 22:35, set by the golden copy, whose latest start is fixed by the features (latest start 23:00) and the backtests (01:00).
\item 22:25: the capture feeds the tick compaction, which feeds the features and so the backtests; the fills feed the positions,
whose latest start (00:40) is set by the risk batch, a later deadline in the chain.
\item The fills: 515 minutes (16:05 to 00:40).
\item Nine: the security master, the positions, the P\&L, the risk report, the reconciliation breaks, the signed P\&L, the
transaction reports, the features and the backtests.
\item The reconciliation breaks and the signed P\&L.
\item 21:55 and 22:15.
\item The security master and the positions (the scenarios and the marks were ready long before).
\item From 22:15 to 22:40, on the new security master and the new positions.
\item 20:15: the file must arrive 45 minutes before the 21:00 start.
\item From 22:15 to 04:15, 165 minutes before the deadline.
\item \textbf{Named result.} On a nominal night the batch ends at 03:55, 185 minutes before 07:00, on the research chain; the
corporate-action file may arrive until 22:35; a risk batch started by the clock at 21:00 reads a stale security master and stale
positions whenever the file arrives after 20:15, only 45 minutes after its nominal time, while started on its inputs it would
have finished at 04:15 on the night of the hook.
\item The stale versions exist and are valid; a fixed-start job has no way to know that a newer version is coming.
\item The third: two versions of the same datasets read at different times.
\item Make the job check that each input's version for the business date is complete before it starts (a completion marker
written by the producer), and refuse or wait otherwise; or let it start but record the version of each input it read, so that
the report is marked stale and rerun when the new version lands.
\item It can be checked and simulated: a program answers what a file reaches, when each report is ready, and which job reads an
incomplete input.
\end{enumerate}
\end{solution}

\begin{solution}{iq:pl:the-platform-map:1}
The one system whose copy of a dataset is authoritative: corrections are made there, and everything else is derived from it.
Without one per dataset, two systems can each be right by their own rules and disagree, and nobody can say which report is
correct.
\iqlookfor{Corrections in one place; copies derived, not edited; disagreement as a symptom of two owners.}
\end{solution}

\begin{solution}{iq:pl:the-platform-map:2}
As a stream when a decision needs it now (pre-trade risk, a trader's screen, a kill switch); as a batch when the answer must be
complete and final (the books, the clearing reconciliation, regulatory reports), because late fills, busts and corporate actions
are only known after the close. Most firms do both and reconcile them.
\iqlookfor{Latency against completeness; both kept, and reconciled.}
\end{solution}

\begin{solution}{iq:pl:the-platform-map:3}
Trace both back through the map to the datasets they read and the versions and times they read them. Check the three causes in
order: two systems of record for one input, a copy made outside the declared flows (a spreadsheet, a cache), or the same dataset
read at two different times (a job that ran before a late input was complete). Then fix the cause, not the number.
\iqlookfor{Lineage and versions rather than guessing; the timing cause.}
\end{solution}

\begin{solution}{iq:pl:the-platform-map:4}
On its inputs, where possible: it starts as early as it can and never reads an incomplete input, but it needs a reliable
completion signal from every producer and can start at an awkward time. A fixed start is predictable for machines and people,
but reads stale data silently when an input is late. If a fixed start is kept, check completeness before starting and record
the versions read.
\iqlookfor{The stale-read failure named; completion markers.}
\end{solution}

\begin{solution}{iq:pl:the-platform-map:5}
A backward pass over the dependency graph: from each report's deadline, subtract job durations back along every path to the
file; the earliest of the resulting latest starts among the jobs that read it is the file's latest safe arrival. Measure the
durations' spread and keep a margin, because the pass assumes fixed durations and free machines.
\iqlookfor{Latest start, backward over the graph, with a margin for variable durations.}
\end{solution}

\begin{solution}{iq:pl:the-platform-map:6}
Start from the outputs people rely on (the P\&L, the risk report, regulatory reports) and trace back what each reads, from
where and when, recording it as data rather than as a drawing; make systems declare their inputs so the map is generated and
checked, not maintained by hand. First checks: one \omterm{def:pl:the-platform-map:sor}{system of record} per dataset, no undeclared copies, and every scheduled job's
inputs complete at its start.
\iqlookfor{Map as data, derived from the systems; ownership and timing checks first.}
\end{solution}
